% Hebrew edition of paper/v2/en.tex in Dnarts/rule110-dependency-fronts.
% Build: xelatex -halt-on-error Rule110-he.tex (run twice for references).
% Fonts: DejaVu Sans, DejaVu Sans Mono, Noto Serif.
\documentclass[11pt,a4paper]{article}
\usepackage[a4paper,margin=2.3cm]{geometry}
\usepackage{fontspec}
\usepackage{amsmath,amssymb,amsthm}
\usepackage[english,bidi=default]{babel}
\babelprovide[import,main]{hebrew}
\babelfont{rm}{Noto Serif}
\babelfont[hebrew]{rm}{DejaVu Sans}
\babelfont{tt}[Scale=0.87]{DejaVu Sans Mono}
\usepackage{booktabs,tabularx}
\usepackage[table]{xcolor}
\usepackage{enumitem}
\usepackage[unicode,hidelinks]{hyperref}
\DeclareRobustCommand{\textenglish}[1]{\foreignlanguage{english}{#1}}
\addto\captionshebrew{%
  \renewcommand{\contentsname}{תוכן העניינים}%
  \renewcommand{\abstractname}{תקציר}%
  \renewcommand{\refname}{ביבליוגרפיה}%
  \renewcommand{\tablename}{טבלה}%
  \renewcommand{\proofname}{הוכחה}%
}
\makeatletter
\def\tagform@#1{\maketag@@@{$\m@th(\mathrm{#1})$}}
\makeatother
\setlength{\emergencystretch}{3em}
\newtheorem{theorem}{משפט}[section]
\newtheorem{lemma}[theorem]{למה}
\newtheorem{proposition}[theorem]{טענה}
\newtheorem{definition}[theorem]{הגדרה}
\newtheorem{conjecture}[theorem]{השערה}
\newtheorem{remark}[theorem]{הערה}
\title{\vspace{-1cm}\textbf{חזיתות תלות מדויקות, תעודות סופיות\texorpdfstring{\\}{} ואינווריאנטים קונסטרוקטיביים לזמן בלתי חסום בכלל 110}\\[0.55em]
{\large מלוא קבוצת הקואורדינטות המהותיות על הסריג האינסופי}}
\author{\textenglish{Aleksandr Strand}}
\date{נובמבר 2026}
\hypersetup{pdftitle={חזיתות תלות מדויקות בכלל 110},pdfauthor={Aleksandr Strand},pdflang={he}}
\begin{document}
\sloppy
\maketitle
\thispagestyle{plain}

\begin{abstract}
אנו קובעים בדיוק אילו קואורדינטות התחלתיות יכולות להשפיע על התא שבראשית הצירים לאחר $t$ צעדים של האוטומט התאי החד־ממדי כלל 110. לכל $t\ge0$, אלה הם כל המספרים השלמים מ־$-(t-\lfloor t/3\rfloor)$ עד $t$, כולל הקצוות. לכל קואורדינטה כזאת קיים זוג תצורות הנבדלות רק בה ומניבות פלט שונה. עבור $t<2485$, אנו מאמתים תעודה סופית על טבעת ומשתמשים במקומיות כדי להעביר את התוצאה לישר האינסופי. עבור $t\ge2485$, חמש רצועות צמודות מכסות את הקטע כולו. באחת הרצועות הפנימיות אוטומט סופי משמר גביש מחזורי עד לשפה נעה, עבור שני הערכים האפשריים של הסיבית הנכנסת. \textenglish{Lean 4.35.0} מאמת את השקילות המתקבלת, \texttt{\textenglish{Rails.essential\_iff}}.
\end{abstract}

\tableofcontents
\vspace{1em}

% ==============================================================================
\section{מבוא והצגת הבעיה}
% ==============================================================================

המשאבים הפיזיים והמידעיים הדרושים לחיזוי מדויק של גודל נצפה בעתיד עשויים להיות שונים במידה ניכרת ממספר הסיביות הדרוש לקידוד מצב אפשרי של המערכת. אנו חוקרים את הגאומטריה של הרגישות המרבית לתנאי ההתחלה בכלל 110, אוטומט תאי אלמנטרי אוניברסלי מבחינה חישובית \cite{cook2004}. הבעיה נשענת על חקר מחלקות שקילות לחיזוי במכניקה חישובית \cite{shalizi2001}, על חקר הגבולות הפיזיקליים של ההיסק \cite{wolpert2008} ועל סיווגם של גליידרים \cite{wolfram1986, martinez2006}. מהירותו של גליידר מתארת את תנועתו של עצם מקומי על רקע מסוים. חסם על הקואורדינטות המהותיות חייב לחול על כל התצורות ההתחלתיות.

\begin{definition}[כלל 110 והגודל הנצפה]
יהי $X = \{0, 1\}^{\mathbb{Z}}$. הצורה האלגברית הנורמלית של הכלל המקומי מעל $\mathbb{F}_2$ היא
\begin{equation}\label{eq:rule110}
f(l, c, r) = c \oplus r \oplus cr \oplus lcr.
\end{equation}
העדכון הסינכרוני $F : X \to X$ נתון על ידי $(F(u))_i = f(u_{i-1}, u_i, u_{i+1})$. הגודל הנצפה בראשית הצירים בזמן $t \in \mathbb{N}$ הוא $G_t(u) = (F^t(u))_0$.
\end{definition}

\begin{definition}[פירוק לצורך חיזוי ומשתנים מהותיים]
מפת תצפית $O : X \to Y$ מספיקה עבור $G_t$ אם קיימת מפה $P : Y \to \{0, 1\}$ כך ש־$G_t = P \circ O$. עבור $j \in \mathbb{Z}$, נסמן את היפוכה של סיבית יחידה באמצעות $(u \oplus e_j)_i = u_i \oplus \mathbf{1}_{i=j}$. קבוצת הקואורדינטות ההתחלתיות המהותיות היא
\begin{equation}\label{eq:essential_set}
D_t = \{ j \in \mathbb{Z} : \exists u \in X,\; G_t(u) \ne G_t(u \oplus e_j) \}.
\end{equation}
\end{definition}

מכיוון ש־$G_t$ תלוי במספר סופי של קואורדינטות בלבד, $D_t$ היא קבוצת קואורדינטות הקלט היחידה המינימלית ביחס להכלה שמספיקה לחישוב $G_t$. כל קטע תצפית מספיק חייב להכיל את הקמור השלם של $D_t$.

כדי לבנות עדויות על רקעים מחזוריים, אנו מבחינים בין קואורדינטת הקלט $j$ לבין הקואורדינטה הפיזית $s$ שבה נצפה פגם. עבור מילה מחזורית $w$, מופע מרחבי $\phi$ והתצורה האינסופית המתאימה $B_0^{w,\phi}$, נגדיר
\begin{equation}\label{eq:diff_support}
B_t^{w,\phi} = F^t(B_0^{w,\phi}), \qquad
Y_t^{w,\phi} = F^t(B_0^{w,\phi} \oplus e_0), \qquad
S_t^{w,\phi} = \{ s \in \mathbb{Z} : B_t^{w,\phi}(s) \ne Y_t^{w,\phi}(s) \}.
\end{equation}
העדכון מתחלף עם הזזות: $F \circ \sigma^d = \sigma^d \circ F$, כאשר $(\sigma^d u)_i = u_{i-d}$. לכל $s \in S_t^{w,\phi}$, הזזת הזוג ההתחלתי ב־$-s$ נותנת עדות לכך ש־$j = -s \in D_t$. מכאן שלכל משפחה $\mathcal{W}$ של רקעים ומופעים מתקיים
\begin{equation}\label{eq:coverage_bridge}
C_t(\mathcal{W}) = \bigcup_{(w,\phi) \in \mathcal{W}} \{ -s : s \in S_t^{w,\phi} \} \subseteq D_t.
\end{equation}
לפיכך, הקצה השמאלי $j = -L(t)$ של חלון הקלט מתאים לחזית הפיזית הימנית $s = L(t)$ של הפגם. הקצה הימני $j = t$ מתאים לחזית הפיזית השמאלית $s = -t$.

% ==============================================================================
\section{חלון התלות המדויק על הסריג האינסופי}
% ==============================================================================

נסמן $L(t) = t - \lfloor t/3 \rfloor$ ו־$I_t = [-L(t), t] \cap \mathbb{Z}$.

\begin{lemma}[חסם תלת־צעדי על התפשטות ימינה]\label{lem:three_step}
עבור מילה בת שבע סיביות $x = (x_0, \dots, x_6) \in \{0, 1\}^7$, יהי $Q(x)$ הערך בתא המרכזי שלה לאחר שלושה צעדים סינכרוניים של כלל 110. אזי
\begin{equation}\label{eq:seven_bit}
Q(x_0, x_1, \dots, x_6) = Q(1 \oplus x_0, x_1, \dots, x_6).
\end{equation}
\end{lemma}

\begin{proof}
לפי \eqref{eq:rule110}, הנגזרת הבוליאנית לפי הארגומנט השמאלי היא
\begin{equation}\label{eq:bool_deriv}
f(1, c, r) \oplus f(0, c, r) = cr.
\end{equation}
נניח ששתי תצורות זהות בכל התאים שמימין ל־$p$, ושההבדל הימני ביותר ביניהן מתקדם תא אחד בכל אחד משלושה צעדים, מ־$p$ ל־$p+1$, משם ל־$p+2$ ולבסוף ל־$p+3$. ההתקדמות הראשונה מחייבת שהערכים המשותפים בתאים $p+1$ ו־$p+2$ יהיו $1$ בזמן $0$. ההתקדמות השנייה מחייבת שהערכים בתאים $p+2$ ו־$p+3$ יהיו $1$ בזמן $1$. מכיוון ש־$f(1, 1, z) = 1 \oplus z$, בהכרח $x_{p+3}(0) = 0$. מכאן, לפי $f(1, 0, z) = z$, בהכרח $x_{p+4}(0) = 1$. לכן $x_{p+4}(1) = f(0, 1, x_{p+5}(0)) = 1$, ומתקבל $x_{p+3}(2) = f(1, 1, 1) = 0$. מכאן שההתקדמות השלישית אינה אפשרית. גם ליבת \textenglish{Lean 4} מאמתת את \eqref{eq:seven_bit} בבדיקת כל $2^7 = 128$ הקלטים.
\end{proof}

\begin{theorem}[חסם אוניברסלי על קבוצת התלות]\label{thm:upper_bound}
לכל $t \in \mathbb{N}$ מתקיים $D_t \subseteq I_t = [-(t - \lfloor t/3 \rfloor), t] \cap \mathbb{Z}$.
\end{theorem}

\begin{proof}
אם שתי תצורות זהות בכל התאים שמימין ל־$p_0$, לאחר צעד אחד הן עדיין זהות בכל התאים שמימין ל־$p_0 + 1$, ולפי למה~\ref{lem:three_step}, לאחר שלושה צעדים הן זהות בכל התאים שמימין ל־$p_0 + 2$. נכתוב $t = 3q + r$, כאשר $q = \lfloor t/3 \rfloor$ ו־$r \in \{0, 1, 2\}$. הזהות נשמרת בכל התאים שמימין ל־$p_0 + 2q + r = p_0 + L(t)$. לכן הבדל ב־$j < -L(t)$ אינו יכול להגיע לראשית הצירים עד הזמן $t$. עבור $j > t$, החסם נובע מחרוט האור הרגיל של כלל בעל רדיוס אחד.
\end{proof}

\begin{theorem}[חדות החסם בשני הקצוות]\label{thm:sharpness}
לכל $t \in \mathbb{N}$, שני הקצוות $-L(t)$ ו־$t$ שייכים ל־$D_t$. קטע התצפית המספיק הקטן ביותר הוא $I_t$. רוחבו מקיים
\begin{equation}
m_{\min}(t) = 2t + 1 - \lfloor t/3 \rfloor, \qquad \lim_{t \to \infty} \frac{m_{\min}(t)}{t} = \frac{5}{3}.
\end{equation}
\end{theorem}

\begin{proof}
עבור הקצה הימני $j = t$, נשווה את קונפיגורציית האפסים $0^{\mathbb{Z}}$ לתצורה שבה יש סיבית אחת שערכה אחד, במיקום $t$, ונסמן אותה ב־$e_t$. כיוון ש־$f(0, 0, 0) = 0$ ו־$f(0, 0, 1) = 1$, הסיבית השמאלית ביותר שערכה אחד נעה בדיוק תא אחד שמאלה בכל צעד ומגיעה לראשית הצירים לאחר $t$ צעדים.

עבור הקצה השמאלי $j = -L(t)$, ניקח את הרקע בעל המחזור ארבע, $B_0(i) = (1, 0, 1, 1)[i \bmod 4]$, ואת ההפרעה שלו $Y_0 = B_0 \oplus e_2$. טבלה~\ref{tab:period4_defect} מציגה את תומך ההפרש $S_t$, המתקבל מהפעלת כלל 110 במשך שישה צעדים.

\begin{table}[h]
\centering
\caption{פגם יחיד על הרקע $(1011)$. לאחר שלב המעבר ההתחלתי, זוג התצורות חוזר על עצמו בהעתק $(\Delta t, \Delta s) = (3, +2)$.}
\label{tab:period4_defect}
\begin{tabular}{cccc}
\toprule
$t$ & מילת הרקע $B_t$ (מחזור $4$) & תומך ההפרש $S_t$ & $\max S_t$ \\
\midrule
$0$ & $1011$ & $\{2\}$ & $2$ \\
$1$ & $1110$ & $\{1, 3\}$ & $3$ \\
$2$ & $1011$ & $\{1, 3, 4\}$ & $4$ \\
$3$ & $1110$ & $\{0, 1, 4\}$ & $4$ \\
$4$ & $1011$ & $\{-1, 0, 1, 3, 5\}$ & $5$ \\
$5$ & $1110$ & $\{3, 5, 6\}$ & $6$ \\
$6$ & $1011$ & $\{2, 3, 6\}$ & $6$ \\
\bottomrule
\end{tabular}
\end{table}

מן המקומיות נובע שהשורות ב־$t=3$ וב־$t=6$ זהות עד כדי הזזה של שני תאים ימינה, ומתקבלות הזהויות $B_6 = \sigma^2 B_3$ ו־$Y_6 = \sigma^2 Y_3$. העובדה שהעדכון מתחלף עם הזזות נותנת $(B_{3+3k+r}, Y_{3+3k+r}) = \sigma^{2k}(B_{3+r}, Y_{3+r})$ לכל $k \ge 0$ ו־$r \in \{0, 1, 2\}$. בצירוף השורות ההתחלתיות, מתקבל $\max S_t = 2 + L(t)$ לכל $t \ge 0$. הזזה של הזוג ההתחלתי ב־$-[2 + L(t)]$ מעבירה את היפוך הסיבית ההתחלתי אל $-L(t)$ ואת ההבדל בזמן $t$ אל ראשית הצירים. לכן $-L(t) \in D_t$.
\end{proof}
\section{טבעות סופיות, קיבולת המצבים וגישה לקלט}
% ==============================================================================

על הטבעת הסופית $\mathbb{Z}_N$ ($N \ge 3$), יהי $X_N = \{0, 1\}^{\mathbb{Z}_N}$ מרחב התצורות ותהי $F_N : X_N \to X_N$ העתקת העדכון. קבוצת המצבים בני־ההשגה בזמן $t$ היא $B_t^{(N)} = F_N^t(X_N)$, ומספר הסיביות הדרושות לקידוד מצב בר־השגה הוא $C_t^{(N)} = \lceil \log_2 |B_t^{(N)}| \rceil$. נגדיר את הדינמיקה של זוגות התצורות $H_N(u, v) = (F_N(u), F_N(v))$, את הזוגות ההתחלתיים $E_j = \{ (u, u \oplus e_j) : u \in X_N \}$ ואת הזוגות הנבדלים בראשית, $W_0 = \{ (x, y) : x_0 \ne y_0 \}$. אזי
\begin{equation}
j \in D_t^{(N)} \iff H_N^t(E_j) \cap W_0 \ne \emptyset.
\end{equation}

\begin{proposition}[תעודה סופית למהותיות מתמשכת]\label{prop:ring_cert}
יהי $\mu_N$ אורכו המרבי של שלב המעבר במערכת הדטרמיניסטית הסופית $F_N$, ויהי $P_N$ הכפולה המשותפת המזערית של אורכי המחזורים שלה. אזי $D_{t+P_N}^{(N)} = D_t^{(N)}$ לכל $t \ge \mu_N$. לכן, לכל $T \in \mathbb{N}$, התנאי $\forall t \ge T,\; D_t^{(N)} = \mathbb{Z}_N$ שקול לבדיקה כי $H_N^t(E_j) \cap W_0 \ne \emptyset$ לכל $j \in \mathbb{Z}_N$ ולכל $T \le t < \max(T, \mu_N) + P_N$.
\end{proposition}

\begin{remark}[דחיסת מצבים ורגישות לקלט]
רגישות בסריג האינסופי אינה נשמרת בהכרח בטבעת סופית. עבור כלל 110 עם $N = 5$, כל $32$ המצבים מגיעים ל־$0^5$ עד לזמן $t = 5$, ולכן $D_t^{(5)} = \emptyset$ לכל $t \ge 5$. בטבעת עם $N = 20$, מנייה ממצה בזמן $t = 121$ נותנת $|B_{121}^{(20)}| = 2095$ מצבים בני־השגה, שאפשר לקודד ב־$12$ סיביות. באותו זמן מתקיים $D_{121}^{(20)} = \mathbb{Z}_{20}$, ולכן כל $20$ סיביות הקלט מהותיות. דוגמאות אלה מבחינות בין ארבעה היבטים של חיזוי: \emph{סיבתיות} (התומך $D_t$), \emph{גישה} (קריאת קואורדינטות בודדות או שימוש בנתונים שעובדו מראש), \emph{קיבולת} ($C_t^{(N)}$) ו\emph{סיבוכיות חישובית}.
\end{remark}

\begin{theorem}[כל קואורדינטה בקטע מהותית]\label{conj:gap_free}
לכל $t \in \mathbb{N}$ מתקיים $D_t = I_t = [-(t - \lfloor t/3 \rfloor), t] \cap \mathbb{Z}$.
\end{theorem}

הבדיקות הסופיות הראשונות כיסו את $0\le t\le2000$. ניסוי מאוחר יותר, שבו השתמשנו במופעים קבועים של שבע מילים, הרחיב את הבדיקה עד $t=5000$. בסעיף האחרון מוכח המשפט לכל הזמנים באמצעות תעודה סופית חדשה עד $t=2484$.

% ==============================================================================
\section{אריחים אוטונומיים חד־צדדיים ושתי רצועות שפה שאינן חסומות בזמן}
% ==============================================================================

מיישום חוזר של למה~\ref{lem:three_step} נובע כי קבוצת התלות של הפלט בתא $s$ אחרי $3m$ צעדים מקיימת $\mathrm{Dep}_{3m}(s) \subseteq [s - 2m, s + 3m]$ לכל התצורות.

\begin{lemma}[סגירות על חצי־ישר]\label{lem:one_sided}
יהי $A_t = F^t(A_0)$, יהי $p > 0$, ונניח כי $a(t + p) = a(t) + \delta$. נניח שלכל זמן בסיס $t_0 \le t < t_0 + p$ מתקיים $A_{t+p}(s + \delta) = A_t(s)$ לכל $s \ge a(t)$. עוד נניח שההגבלה של $A_{t+p}$ ל־$[a(t+p), \infty)$ תלויה רק בהגבלה של $A_t$ ל־$[a(t), \infty)$ בכל צעד באורך $p$. אזי $A_{t+p}(s + \delta) = A_t(s)$ לכל $t \ge t_0$ ולכל $s \ge a(t)$. הטענה המקבילה מתקיימת גם בחצאי־ישר שמאליים $s \le a(t)$.
\end{lemma}

ניישם את למה~\ref{lem:one_sided} על החזית הימנית במרחב הפיזי $e(t) = L(t)$ עם $(\Delta t, \Delta s) = (30, +20)$. כאן $e(t+30) = e(t) + 20$, וקבוצת התלות של האריח החדש שרוחבו $W$ מוכללת בקטע $[e(t) - W + 1, e(t) + 50]$. קצהו השמאלי של הקטע מותיר מחוצה לו את התחום הפנימי הלא־ידוע $s < e(t) - W + 1$. יישום אותו נימוק על החזית השמאלית $s = -t$ עם $(\Delta t, \Delta s) = (40,-40)$ נותן את התוצאה הבאה, ללא חסם עליון על הזמן.

\begin{theorem}[שתי רצועות שפה]\label{thm:dual_strips}
לכל מספר שלם $t \ge 300$,
\begin{equation}\label{eq:dual_strips}
[-L(t), -L(t) + 105] \cap \mathbb{Z} \;\cup\; [t - 7, t] \cap \mathbb{Z} \;\subseteq\; D_t.
\end{equation}
\end{theorem}

התעודה לרצועה השמאלית, שרוחבה $106$, משתמשת ב־$20$ מופעים של הרקעים $(0001111001)$, $(1011)$ ו־$(001011)$, עם עומק אריח של $167$ תאים, על פני $30$ זמני הבסיס $t = 300, \dots, 329$. התעודה לרצועה הימנית, שרוחבה $8$, משתמשת בשישה מופעים עם $(\Delta t, \Delta s) = (40,-40)$ על פני זמני הבסיס $t = 300, \dots, 339$. \emph{חישוב מלווה בהוכחה} על וקטורי סיביות מטיפוס \textenglish{\texttt{BitVec 1320}} מאמת את מצבי הביניים. הליבה של \textenglish{Lean 4.35.0} בודקת כל מעבר $F_{1320}^{h_i}(z_i) = z_{i+1}$ עבור $h_i \in \{1, 10\}$.

% ==============================================================================
\section{קרניים בליסטיות ומגבלותיהן של תעודות מקומיות}
% ==============================================================================

דרך נוספת להרחיב את הכיסוי לזמנים גדולים כרצוננו היא לשמר הבדל לאורך מסלול בדיד במרחב ובזמן.

\begin{lemma}[התמדה על חצי־ישר ימני]\label{lem:ballistic_ray}
נניח שמסלול $A$ מקיים $\forall s \ge x,\; A_{T+3}(s + 2) = A_T(s)$ בזמן $T$. אזי $\forall q \ge 0,\; \forall s \ge x,\; A_{T+3q}(s + 2q) = A_T(s)$. בפרט, אם גם $B$ וגם $Y$ מקיימים תנאי זה ומתקיים $B_T(x) \ne Y_T(x)$, אזי $B_{T+3q}(x + 2q) \ne Y_{T+3q}(x + 2q)$ לכל $q \ge 0$.
\end{lemma}

נתבונן בשתי המשפחות החשבוניות של נקודות ההתחלה על הרקעים $(0000101111)$ ו־$(001011)$, עם $0 \le n \le 32$:
\begin{equation}
T_A(n) = 594 + 882n,\quad x_A(n) = 256 + 378n; \qquad
T_B(n) = 884 + 882n,\quad x_B(n) = 380 + 378n,
\end{equation}
בהתאמה. חישוב אימת את $66$ נקודות ההתחלה הללו. כל אחת מהן יוצרת קרן בליסטית אינסופית של קואורדינטות מהותיות בעומק קבוע ביחס לחזית $e(t) = L(t)$. העומקים הם $d_A(n) = 140 + 210n$ ו־$d_B(n) = 210 + 210n$, בהתאמה.

המודולים הבאים, \textenglish{\texttt{StaircaseA}} ו־\textenglish{\texttt{StaircaseB}}, מוכיחים את הנוסחאות לשתי המשפחות לכל $n,q\ge0$, באמצעות חיבור של שלוש שורות התואמות זו לזו מקומית לכדי מסלול מדויק. התווכים המחזוריים בחזיתות של המשפחות \textenglish{A} ו־\textenglish{B} מכסים $45\,255$ מתוך $66\,150$ מחלקות השאריות. הוספת משפחה שלישית מעלה את המספר ל־$53\,550$. מניינים אלה מתייחסים לטבלת השאריות שממדיה $315\times210$. כדי להחיל תא בטבלה על קואורדינטה במרחב הפיזי, יש לבדוק גם שהנקודה שייכת לטריז המתאים.

דוגמאות נגדיות מגלות שלוש מגבלות של תעודות אלה.
\begin{itemize}
    \item \textbf{אריחים מקומיים אינם מספיקים למעבר האינדוקציה מ־$n$ ל־$n+1$.} עבור הצעד $(\Delta T, \Delta x) = (882, +378)$, חסם התלות האוניברסלי הוא $\mathrm{Dep}_{882}(s) \subseteq [s - 588, s + 882]$. בנקודת היעד לאחר ההזזה מתקבל הקטע $[x - 210, x + 1260]$. הדוגמאות הנגדיות שבקובץ \textenglish{\texttt{check\_macro\_boundary.py}} מקבעות אריחים ברדיוס $64$ (ב־$T=594$) או $186$ (ב־$T=884$). בכל מקרה, שינוי תאים מחוץ לאריח משנה את סיבית היעד אחרי $882$ צעדים.
    \item \textbf{בגזרה שנצפתה לאורך זמן סופי עשויים להתגלות פערים בהמשך.} ברקע $(0001111001)$, הגזרה השמאלית שבה המהירויות נמצאות בקטע $[-0.66, -0.64]$ נראית מכוסה עד $t = 2000$. הרצה ארוכה יותר חושפת פער.
    \item \textbf{חזרה עד כדי הזזה אינה מספיקה כשלעצמה להוכחת כיסוי.} אפילו חזרה עד לעומק של $261$ תאים עבור $(0000001001)$ אינה מבטיחה כיסוי של כל מחלקות השאריות. הוכחת הכיסוי משלבת כמה מופעים.
\end{itemize}

\section{השלמה קונסטרוקטיבית של הגזרה השמאלית הבלתי חסומה}
% ==============================================================================

כדי להביא בחשבון את התנועות השונות של השפה השמאלית של חרוט האור (שמהירותה $-1$), של הגביש בעל המחזור $50$ ושל חזית ההפרעה, נאנדקס את הזמנים הזוגיים באמצעות $j \in \mathbb{N}$ ונגדיר
\begin{equation}
\mathrm{tm}(j) = 2450 + 2j, \qquad a(j) = -350 - j.
\end{equation}
הגבול הימני $a(j)$ נע תא אחד שמאלה בכל שני צעדים, במהירות $-1/2$. לכל אחת משש המופעים $p \in \{\mathrm{p0}, \mathrm{p1}, \mathrm{p2}, \mathrm{p3}, \mathrm{p7}, \mathrm{p9}\}$, התנאי האינווריאנטי בשכבה $j$ מורכב משני חלקים. המסלול תואם לגביש הייחוס $\mathrm{Cm}_p$ בקטע שמשמאל ל־$a(j)$, ומאגר הגבול שלו, בן $16$ הסיביות, מכיל את המילה השמורה בקודקוד בגרף מעברים סופי:
\begin{align}
\mathrm{Medium}(p, j) &\iff \forall x \in \mathbb{Z},\quad -\mathrm{tm}(j) + 1300 \le x < a(j) \implies Y_p(\mathrm{tm}(j), x) = \mathrm{Cm}_p(\mathrm{tm}(j), x), \\
\mathrm{Window}(p, j, i) &\iff \forall k \in \{0, \dots, 15\},\quad Y_p(\mathrm{tm}(j), a(j) + k) = \mathrm{testBit}(\mathrm{word}(p, j, i), k).
\end{align}

\begin{lemma}[מתמר דטרמיניסטי בן 16 סיביות והממשק שלו עם הרצועה השמאלית]\label{lem:transducer}
לכל מופע $p$, הטבלה \textenglish{\texttt{Crystal50Packed}} מגדירה גרף ששכבותיו חוזרות במחזור $50$. כל רשומה תופסת $36$ סיביות. $16$ הסיביות הפחות משמעותיות מקודדות את החלון $\mathrm{word}(p, j, i)$, והסיביות $16\dots 25$ ו־$26\dots 35$ מקודדות את המעברים עבור שני הערכים האפשריים של סיבית הקלט הימנית, $r \in \{0, 1\}$.
\begin{enumerate}[label=(\alph*)]
    \item בצעד $j$, ההקשר השמאלי בן $4$ הסיביות $\mathrm{ctx}(p, j \bmod 50)$ תופס את המקומות $a(j) - 4 + m$ עבור $m \in \{0, 1, 2, 3\}$. עומקם הוא $t + x = 2096 + j + m$, והם נמצאים בקטע שכבר מכוסה על ידי $\mathrm{Medium}(p, j)$.
    \item הפעלה כפולה של פעולת הצמצום הסיבית, $\mathrm{shrink}(\mathrm{shrink}(\mathrm{input}(p, j, i, r)))$, מחשבת את $17$ סיביות התצורה בקטע $[a(j)-2, a(j)+14]$ בזמן $\mathrm{tm}(j+1)$. הסיבית הפחות משמעותית ($k=0$) נותנת את $\mathrm{Cm}_p(\mathrm{tm}(j+1), a(j) - 2)$. יתר $16$ הסיביות ($k=1\dots 16$) יוצרות את המילה בקודקוד $z$ בשכבה הבאה, שחלונה מתחיל ב־$a(j+1) = a(j) - 1$.
    \item בשפה השמאלית של הרצועה, עומקי התאים $-\mathrm{tm}(j+1) + 1300 \le x < -\mathrm{tm}(j) + 1302$ הם $k = x + \mathrm{tm}(j+1) \in \{1300,1301,1302,1303\}$. המשפט \textenglish{\texttt{seam}} במודול \textenglish{\texttt{Crystal50Seam}} מכסה תאים אלה, שכן $t + x$ אינו משתנה תחת ההעתקה $(t, x) \mapsto (t + 50, x - 50)$.
\end{enumerate}
\end{lemma}

\begin{theorem}[\textenglish{\texttt{Rails.Crystal50Wedge.wedge\_all}}]\label{thm:wedge_all}
לכל זמן שלם $t \ge 2485$ ולכל $x$ שלם המקיים
\begin{equation}\label{eq:wedge_region}
-t + 1302 \le x \le 59 - \lfloor t/2 \rfloor,
\end{equation}
מתקיים הפרדיקט \textenglish{\texttt{Rule110KernelCert.EssentialAt t (-x)}}. אין חסם עליון על הזמן $t$.
\end{theorem}

\begin{proof}[מתווה ההוכחה]
נגדיר $j = \lfloor (t - 2450) / 2 \rfloor$. עבור $t \ge 2485$ מתקיים $j \ge 17$, והגבול הימני הוא $59 - \lfloor t/2 \rfloor = -1166 - j = a(j) - 816$. בזמנים זוגיים, $t = \mathrm{tm}(j)$, התנאי האינווריאנטי $\mathrm{Medium}(p, j)$ נותן ישירות $Y_p(t, x) = \mathrm{Cm}_p(t, x)$ לכל שש המופעים $p$. בזמנים אי־זוגיים, $t = \mathrm{tm}(j) + 1$, קטע התלות של צעד אחד $[x - 1, x + 1]$ מוכל כולו ב־$[-\mathrm{tm}(j) + 1300, a(j))$. לכן למת המקומיות \textenglish{\texttt{carry}} מעבירה את ההתאמה לצעד הבא. מהתאמה זו מתקבלת עדות למהותיות באמצעות \textenglish{\texttt{essential\_of\_cm}}.
\end{proof}

% ==============================================================================
\section{כיסוי מלא על הישר האינסופי}
נשתמש בקואורדינטת הקלט $j=-x$ ונכתוב $b=\lfloor t/2\rfloor$. עבור $t\ge2485$, חמש התוצאות הבאות חלות על קטעים עוקבים. בכל שורה שני הקצוות כלולים.
\begin{center}
\begin{tabular}{lll}
\toprule
קואורדינטה ראשונה & קואורדינטה אחרונה & תוצאה \\
\midrule
$-L(t)$ & $-L(t)+105$ & תעודת חצי־ישר ברוחב 106 \\
$-L(t)+106$ & $b-60$ & \textenglish{\texttt{EtherWedge.half\_sector}} \\
$b-59$ & $t-1302$ & \textenglish{\texttt{Crystal50Wedge.wedge\_all}} \\
$t-1301$ & $t-20$ & \textenglish{\texttt{LeftCrystal50.left\_strip}} \\
$t-19$ & $t$ & \textenglish{\texttt{right\_boundary\_width20}} \\
\bottomrule
\end{tabular}
\end{center}
שלוש מופעים של האתר בעל המחזור 14 מבססות את השורה השנייה לכל $t\ge210$. השורה האחרונה מרחיבה את הרצועה הימנית הקודמת, שרוחבה שמונה. בזמן $t=2485$, השורה השלישית מכילה תא אחד בלבד, $j=1183$. כל שורה מתחילה בקואורדינטה העוקבת לקואורדינטה האחרונה של השורה הקודמת. הלמה \textenglish{\texttt{Rails.late\_all}} משלבת את חמש ההוכחות.

\paragraph{התעודה הסופית.}
עבור $0\le t<2485$, נשתמש בכל המופעים המחזוריות של שבע מילים. הזוגות שלהלן מציינים את אורכה ואת הקידוד השלם של כל מילה, כאשר התא הראשון מיוצג בסיבית הפחות משמעותית:
\[
(8,244),\ (10,632),\ (10,976),\ (10,576),\ (4,13),\ (6,52),\ (6,48).
\]
אנו מייצגים את הרקעים כווקטורים של סיביות באורך 5520, והופכים סיבית בודדת במקום 2760. המחזורים של כל שבע המילים מחלקים את 5520. בכל צעד זמן אנו מחשבים את האיחוד הסיבי של ההפרשים על פני כל הזוגות ומוודאים שהוא מכיל את המסכה של הקטע הפיזי $[-t,L(t)]$. כך מכוסים $\sum_{t=0}^{2484}(2t+1-\lfloor t/3\rfloor)=5\,147\,263$ זוגות מרחב־זמן. לכל $x$ בקטע זה מתקיים $-5520<x-t$ וכן $x+t<5520$. לכן למת המקומיות \textenglish{\texttt{local\_obs}} מעבירה את התוצאה להיפוך סיבית בודדת על הישר האינסופי.

שבע המילים קובעות את השורה ההתחלתית. המחולל מספק תמונות מצב מועמדות בכל 25 צעדים. \textenglish{Lean} בודקת את מסכות הכיסוי של כל בלוק ומוודאת שמצבו הסופי שווה לתמונת המצב ההתחלתית של הבלוק הבא, אם יש כזה. הבלוק האחרון מכיל עשר שורות. אינדוקציה על אורך הבלוק משייכת לכל סיבית מכוסה במסכה זוג מסלולים מסוים. בדרך זו \textenglish{\texttt{Rails.Finite.finite\_all}} מוכיחה מהותיות לכל קואורדינטת קלט בתחום הסופי.

\begin{proof}[השלמת הוכחת משפט~\ref{conj:gap_free}]
ממשפט~\ref{thm:upper_bound} נובע $D_t\subseteq I_t$. להכלה ההפוכה, נפעיל את התעודה הסופית עבור $t<2485$ ואת חמש הרצועות עבור $t\ge2485$. הזזה של כל זוג הנבדל בקואורדינטה הפיזית $x$ נותנת עדות לכך שקואורדינטת הקלט $-x$ מהותית.
\end{proof}

\paragraph{מסקנה בנוגע לגישה לנתונים.}
כל קבוצה קבועה של קואורדינטות התחלתיות המספיקה לחישוב מדויק של $G_t$ בכל התצורות מכילה את $I_t$. מן המקומיות והחסם העליון נובע ש־$I_t$ עצמה מספיקה. המספר המזערי של קואורדינטות בתצפית כזו הוא אפוא $2t+1-\lfloor t/3\rfloor$. העדות יכולה להיות תלויה הן ב־$t$ והן ב־$j$. טענות על עלות של שאילתות אדפטיביות, על עיבוד מקדים או על זיכרון מחייבות מודל חישובי נפרד.

\paragraph{שחזור התוצאות.}
משפט השורש נמצא בקובץ \textenglish{\texttt{lean/Rule110Exact.lean}}. הקוד הנלווה כולל את קובצי המקור של המודולים המיובאים, את מחולל הבלוקים הסופיים, את סקריפט הבנייה ואת יומני הבנייה. הפקודה \textenglish{\texttt{\#print axioms}} מחזירה \textenglish{\texttt{[propext, Classical.choice, Quot.sound]}} עבור משפט השורש, ו־\textenglish{\texttt{[propext, Quot.sound]}} עבור \textenglish{\texttt{wedge\_all}}. כל רשימה מציינת את התלויות של המשפט המתאים.

\begin{thebibliography}{9}
\bibitem{cook2004}
\textenglish{M.~Cook, ``Universality in elementary cellular automata,'' \emph{Complex Systems}, vol.~15, no.~1, pp.~1--40, 2004.}

\bibitem{martinez2006}
\textenglish{G.~J.~Mart\'inez, H.~V.~McIntosh, and J.~C.~Seck-Tuoh-Mora, ``Gliders in Rule 110,'' \emph{International Journal of Unconventional Computing}, vol.~2, no.~1, pp.~1--49, 2006.}

\bibitem{shalizi2001}
\textenglish{C.~R.~Shalizi and J.~P.~Crutchfield, ``Computational mechanics: Pattern and prediction, structure and simplicity,'' \emph{Journal of Statistical Physics}, vol.~104, pp.~817--879, 2001.}

\bibitem{wolfram1986}
\textenglish{S.~Wolfram, ``Tables of cellular automaton properties,'' in \emph{Theory and Applications of Cellular Automata}, World Scientific, pp.~485--557, 1986.}

\bibitem{wolpert2008}
\textenglish{D.~H.~Wolpert, ``Physical limits of inference,'' \emph{Physica D: Nonlinear Phenomena}, vol.~237, no.~9, pp.~1257--1281, 2008.}
\end{thebibliography}

\end{document}
